\apendice{Roteiro de reprodução do estudo computacional}
\label{ap:roteiro-reproducao}

Este apêndice apresenta uma sequência de operações para reproduzir, na ferramenta desenvolvida, a análise e o pré-dimensionamento do modelo descrito no Capítulo~\ref{chap:resultados}. O roteiro se refere à avaliação bidimensional de uma face da ponte e não constitui um procedimento de ensaio físico.

\begin{enumerate}
    \item Preparar no FTOOL o modelo correspondente a uma face da ponte e salvá-lo como \texttt{ponte.ftl}, preservando a geometria, os apoios, as propriedades estruturais e o carregamento vertical de \(2\,\mathrm{kN}\) utilizados no estudo.
    \item Abrir a aplicação web e selecionar o arquivo \texttt{ponte.ftl} para importação.
    \item Conferir a prévia da geometria reconstruída, verificando a identificação de 10 nós e 19 barras.
    \item Conferir as propriedades e os parâmetros de pré-dimensionamento adotados no estudo, incluindo o comprimento comercial de \(115\,\mathrm{mm}\), a sobreposição padrão de \(0\,\mathrm{mm}\) e o mínimo construtivo de três camadas por barra. Os demais parâmetros devem corresponder aos adotados na execução de referência.
    \item Executar a análise estrutural e examinar o diagrama de esforços axiais e os resultados numéricos das barras.
    \item Executar o pré-dimensionamento e consultar a distribuição de camadas e a estimativa de palitos comerciais por barra.
    \item Consultar o relatório consolidado e confrontar seus resultados com as tabelas e figuras apresentadas no Capítulo~\ref{chap:resultados}.
\end{enumerate}

Na configuração utilizada para os resultados de referência, a ferramenta apresentou uma soma de 154 camadas nas 19 barras e uma estimativa de 430 palitos comerciais para a face analisada. A extrapolação para 860 palitos considera duas faces idênticas, sem incluir elementos de ligação transversal ou perdas de fabricação.

Esses valores servem como referência para a reprodução do caso estudado e dependem da manutenção das propriedades do modelo estrutural e dos parâmetros de pré-dimensionamento utilizados na avaliação.
